\chapter{Preprocessing}

\section{Introduction}

Before a C compiler sees your source code, it passes through the \textbf{preprocessor}---a powerful text transformation engine that manipulates your code at the character level. The preprocessor doesn't understand C syntax; it simply performs textual substitutions, file inclusions, and conditional deletions.

Understanding preprocessing is essential because:
\begin{itemize}
    \item It's where macros, includes, and conditionals are handled
    \item Errors in preprocessing can produce baffling compiler errors
    \item The preprocessor determines what the compiler actually sees
    \item Many build configurations rely on preprocessor conditionals
\end{itemize}

In this chapter, we'll explore every aspect of the C preprocessor, from simple macros to complex conditional compilation, with hands-on examples you can try yourself.

\section{What is the Preprocessor?}

The C preprocessor (often abbreviated as CPP) is a separate program that processes your source code before the compiler proper. On most systems, it's invoked automatically by the compiler driver (like \texttt{gcc}), but you can also run it standalone:

\begin{lstlisting}[language=bash]
# Run preprocessor only, output to stdout
gcc -E source.c

# Run preprocessor only, output to file
gcc -E source.c -o source.i
\end{lstlisting}

The preprocessor operates on \textbf{translation units}---essentially, what the compiler processes as a single entity. The output of the preprocessor is a single, expanded text stream that the compiler then parses.

\subsection{Preprocessing as Text Transformation}

The preprocessor performs three main operations:

\begin{enumerate}
    \item \textbf{Macro expansion}: Replace defined symbols with their definitions
    \item \textbf{File inclusion}: Insert the contents of other files
    \item \textbf{Conditional compilation}: Include or exclude code based on conditions
\end{enumerate}

All of these are pure text operations. The preprocessor doesn't know about C syntax, types, or scope. It just manipulates text.

\begin{verbatim}
┌─────────────────────────────────────────────────────────────┐
│              The Preprocessing Pipeline                     │
├─────────────────────────────────────────────────────────────┤
│                                                             │
│   source.c ──┐                                              │
│              │    ┌─────────────┐    ┌─────────────┐        │
│   header.h ──┼───▶│ Preprocessor│───▶│ source.i    │        │
│              │    │   (cpp)     │    │ (expanded)  │        │
│   config.h ──┘    └─────────────┘    └─────────────┘        │
│                         │                                   │
│                    Text transformations:                    │
│                    - #include → file contents               │
│                    - #define → substitutions                │
│                    - #if/#ifdef → conditional code          │
│                                                             │
└─────────────────────────────────────────────────────────────┘
\end{verbatim}

\section{Macro Expansion}

\subsection{Object-Like Macros}

The simplest macros are object-like---they replace a name with a token sequence:

\begin{lstlisting}[language=C]
#define PI 3.14159
#define MAX_SIZE 1024
#define DEBUG 1

int main(void) {
    double area = PI * radius * radius;
    int buffer[MAX_SIZE];
    if (DEBUG) {
        printf("Debug mode enabled\n");
    }
    return 0;
}
\end{lstlisting}

After preprocessing, this becomes:

\begin{lstlisting}[language=C]
int main(void) {
    double area = 3.14159 * radius * radius;
    int buffer[1024];
    if (1) {
        printf("Debug mode enabled\n");
    }
    return 0;
}
\end{lstlisting}

\begin{tip}
Use UPPERCASE names for macros to distinguish them from variables. This is a convention, not a requirement, but it helps readers identify macros at a glance.
\end{tip}

\subsection{Function-Like Macros}

Macros can take arguments, making them look like functions:

\begin{lstlisting}[language=C]
#define MAX(a, b) ((a) > (b) ? (a) : (b))
#define SQUARE(x) ((x) * (x))
#define ABS(x) ((x) < 0 ? -(x) : (x))

int main(void) {
    int m = MAX(3, 5);          // Expands to: ((3) > (5) ? (3) : (5))
    int sq = SQUARE(4);          // Expands to: ((4) * (4))
    int av = ABS(-10);           // Expands to: ((-10) < 0 ? -(-10) : (-10))
    return 0;
}
\end{lstlisting}

\begin{warning}
\textbf{Critical Warning}: Always parenthesize macro arguments! Without parentheses, operator precedence can cause unexpected behavior:

\begin{lstlisting}[language=C]
#define BAD_SQUARE(x) x * x
int result = BAD_SQUARE(1 + 2);  // Expands to: 1 + 2 * 1 + 2 = 5, not 9!
\end{lstlisting}
\end{warning}

\subsection{Macro Pitfall: Multiple Evaluation}

Function-like macros can evaluate arguments multiple times:

\begin{lstlisting}[language=C]
#define MAX(a, b) ((a) > (b) ? (a) : (b))

int x = 5;
int m = MAX(x++, 3);  // Problem: x might be incremented once or twice!
\end{lstlisting}

After expansion:
\begin{lstlisting}[language=C]
int m = ((x++) > (3) ? (x++) : (3));
\end{lstlisting}

If \texttt{x} is 5, then \texttt{x++} (which returns 5) is compared to 3. Since 5 > 3 is true, the result is \texttt{x++} again---but now \texttt{x} is 6, so this returns 6. \texttt{x} ends up as 7, and \texttt{m} is 6.

\textbf{Solution}: Use inline functions when possible, or be very careful with macro arguments that have side effects:

\begin{lstlisting}[language=C]
// Prefer inline functions for type safety and single evaluation
static inline int max(int a, int b) {
    return a > b ? a : b;
}
\end{lstlisting}

\subsection{Multi-Line Macros}

For longer macros, use backslash continuation:

\begin{lstlisting}[language=C]
#define SWAP(a, b, type) do { \
    type temp = a; \
    a = b; \
    b = temp; \
} while (0)

int main(void) {
    int x = 1, y = 2;
    SWAP(x, y, int);  // x = 2, y = 1
    return 0;
}
\end{lstlisting}

\begin{tip}
\textbf{Why the \texttt{do \{ ... \} while (0)}?} This idiom allows the macro to be used like a statement with a semicolon, without breaking if-else chains:

\begin{lstlisting}[language=C]
if (condition)
    SWAP(a, b, int);  // Works correctly
else
    do_something();
\end{lstlisting}
\end{tip}

\subsection{Stringification (\#)}

The \texttt{\#} operator converts a macro argument to a string:

\begin{lstlisting}[language=C]
#define STRINGIFY(x) #x
#define TOSTRING(x) STRINGIFY(x)

#define VALUE 42

const char* s1 = STRINGIFY(VALUE);   // s1 = "VALUE"
const char* s2 = TOSTRING(VALUE);    // s2 = "42"
\end{lstlisting}

The difference: \texttt{STRINGIFY} stringifies the argument literally, while \texttt{TOSTRING} first expands the argument, then stringifies it.

\subsection{Token Pasting (\#\#)}

The \texttt{\#\#} operator concatenates tokens:

\begin{lstlisting}[language=C]
#define CONCAT(a, b) a ## b
#define MAKE_VAR(name, num) int name ## _ ## num

int main(void) {
    int CONCAT(var, 1) = 10;        // Creates variable var1
    MAKE_VAR(item, 5) = 20;          // Creates variable item_5
    return 0;
}
\end{lstlisting}

This is commonly used for generating names programmatically:

\begin{lstlisting}[language=C]
#define DEFINE_HANDLER(num) void handler_##num(void) { \
    printf("Handler %d called\n", num); \
}

DEFINE_HANDLER(1)  // Creates handler_1
DEFINE_HANDLER(2)  // Creates handler_2
DEFINE_HANDLER(3)  // Creates handler_3
\end{lstlisting}

\subsection{Variadic Macros}

Macros can accept variable numbers of arguments using \texttt{...} and \texttt{\_\_VA\_ARGS\_\_}:

\begin{lstlisting}[language=C]
#define DEBUG_PRINT(fmt, ...) \
    fprintf(stderr, "[DEBUG] " fmt "\n", __VA_ARGS__)

#define LOG(level, msg, ...) \
    printf("[%s] " msg "\n", level, __VA_ARGS__)

int main(void) {
    DEBUG_PRINT("Value: %d, Status: %s", 42, "OK");
    LOG("INFO", "User %s logged in", "alice");
    return 0;
}
\end{lstlisting}

For C99 compatibility with zero variadic arguments, use \texttt{\#\_\_VA\_ARGS\_\_} (GCC extension, now standard in C++20):

\begin{lstlisting}[language=C]
#define DEBUG_PRINT(fmt, ...) \
    fprintf(stderr, "[DEBUG] " fmt "\n", ##__VA_ARGS__)

DEBUG_PRINT("Simple message");  // Works even with no extra args
\end{lstlisting}

\section{File Inclusion}

\subsection{\#include Basics}

The \texttt{\#include} directive inserts the contents of another file:

\begin{lstlisting}[language=C]
#include <stdio.h>      // System header (angle brackets)
#include "myheader.h"   // Local header (double quotes)
\end{lstlisting}

The difference between \texttt{<...>} and \texttt{"..."} is the search path:

\begin{tabular}{|l|p{10cm}|}
\hline
\textbf{Syntax} & \textbf{Search Order} \\
\hline
\texttt{<header.h>} & System include paths only \\
\hline
\texttt{"header.h"} & Current directory first, then system include paths \\
\hline
\end{tabular}

\subsection{Include Search Paths}

The preprocessor searches for headers in specific directories:

\begin{lstlisting}[language=bash]
# View default include paths
gcc -E -x c - -v < /dev/null 2>&1 | grep "^ "

# Add custom include paths
gcc -I/path/to/includes -I../another/path source.c
\end{lstlisting}

Common system include paths on Linux:
\begin{itemize}
    \item \texttt{/usr/include}
    \item \texttt{/usr/local/include}
    \item \texttt{/usr/lib/gcc/x86\_64-linux-gnu/11/include} (compiler-specific)
\end{itemize}

\subsection{The Include Graph}

With nested includes, you can build complex dependency graphs:

\begin{verbatim}
source.c
├── <stdio.h>
│   ├── <bits/libc-header-start.h>
│   ├── <features.h>
│   │   ├── <stdc-predef.h>
│   │   └── <sys/cdefs.h>
│   └── <bits/types.h>
├── "config.h"
│   └── <stdint.h>
└── "utils.h"
    └── <stdlib.h>
\end{verbatim}

To see the include graph:

\begin{lstlisting}[language=bash]
# Show included files
gcc -E source.c -H

# Output example:
# . /usr/include/stdio.h
# .. /usr/include/bits/libc-header-start.h
# .. /usr/include/features.h
\end{lstlisting}

\subsection{Include Guards and \#pragma once}

Without protection, a header included twice causes errors:

\begin{lstlisting}[language=C]
// header.h
struct Point { int x, y; };

// source.c
#include "header.h"
#include "header.h"  // Error: redefinition of 'struct Point'!
\end{lstlisting}

\textbf{Solution 1: Include Guards (Traditional)}

\begin{lstlisting}[language=C]
#ifndef HEADER_H
#define HEADER_H

struct Point { int x, y; };

#endif /* HEADER_H */
\end{lstlisting}

The first inclusion defines \texttt{HEADER\_H}. Subsequent inclusions see the guard and skip the content.

\textbf{Solution 2: \#pragma once (Modern)}

\begin{lstlisting}[language=C]
#pragma once

struct Point { int x, y; };
\end{lstlisting}

\texttt{\#pragma once} tells the preprocessor to include this file at most once per compilation unit. It's simpler and often faster, but technically non-standard (though supported by all major compilers).

\begin{recommendation}
Use \texttt{\#pragma once} for new code. It's cleaner and less error-prone. Use include guards if you need strict C standard compliance.
\end{recommendation}

\section{Conditional Compilation}

Conditional compilation lets you include or exclude code based on conditions evaluated during preprocessing.

\subsection{\#if, \#elif, \#else, \#endif}

\begin{lstlisting}[language=C]
#define VERSION 3

#if VERSION == 1
    #define FEATURE_NAME "Basic"
#elif VERSION == 2
    #define FEATURE_NAME "Standard"
#elif VERSION == 3
    #define FEATURE_NAME "Premium"
#else
    #define FEATURE_NAME "Unknown"
#endif
\end{lstlisting}

\subsection{\#ifdef and \#ifndef}

Test whether a macro is defined:

\begin{lstlisting}[language=C]
#ifdef DEBUG
    printf("Debug mode enabled\n");
#endif

#ifndef BUFFER_SIZE
    #define BUFFER_SIZE 1024  // Default value
#endif
\end{lstlisting}

These are equivalent to:

\begin{lstlisting}[language=C]
#if defined(DEBUG)
#if !defined(BUFFER_SIZE)
\end{lstlisting}

\subsection{defined() Operator}

The \texttt{defined()} operator checks if a macro exists:

\begin{lstlisting}[language=C]
#if defined(DEBUG) && defined(VERBOSE)
    #define LOG_LEVEL 3
#elif defined(DEBUG)
    #define LOG_LEVEL 2
#else
    #define LOG_LEVEL 1
#endif
\end{lstlisting}

\subsection{Common Use Cases}

\textbf{Platform Detection:}

\begin{lstlisting}[language=C]
#ifdef _WIN32
    #define PATH_SEPARATOR '\\'
    #include <windows.h>
#elif defined(__linux__)
    #define PATH_SEPARATOR '/'
    #include <unistd.h>
#elif defined(__APPLE__)
    #define PATH_SEPARATOR '/'
    #include <mach-o/dyld.h>
#endif
\end{lstlisting}

\textbf{Debug vs. Release:}

\begin{lstlisting}[language=C]
#ifdef NDEBUG
    #define ASSERT(cond) ((void)0)
#else
    #define ASSERT(cond) \
        do { \
            if (!(cond)) { \
                fprintf(stderr, "Assertion failed: %s, %s:%d\n", \
                        #cond, __FILE__, __LINE__); \
                abort(); \
            } \
        } while (0)
#endif
\end{lstlisting}

\textbf{Feature Flags:}

\begin{lstlisting}[language=C]
// Compile with: gcc -DUSE_OPENCL source.c
#ifdef USE_OPENCL
    void compute(void) { /* OpenCL implementation */ }
#elif defined(USE_CUDA)
    void compute(void) { /* CUDA implementation */ }
#else
    void compute(void) { /* CPU fallback */ }
#endif
\end{lstlisting}

\textbf{Header Inclusion Control:}

\begin{lstlisting}[language=C]
// Disable assert.h assertions
#define NDEBUG
#include <assert.h>
\end{lstlisting}

\section{Predefined Macros}

The C standard and compilers provide predefined macros:

\subsection{Standard Predefined Macros}

\begin{tabular}{|l|l|p{8cm}|}
\hline
\textbf{Macro} & \textbf{Description} & \textbf{Example Value} \\
\hline
\texttt{\_\_FILE\_\_} & Current source file name & \texttt{"source.c"} \\
\hline
\texttt{\_\_LINE\_\_} & Current line number & \texttt{42} \\
\hline
\texttt{\_\_func\_\_} & Current function name (C99) & \texttt{"main"} \\
\hline
\texttt{\_\_DATE\_\_} & Compilation date & \texttt{"Mar 11 2026"} \\
\hline
\texttt{\_\_TIME\_\_} & Compilation time & \texttt{"16:30:00"} \\
\hline
\texttt{\_\_STDC\_\_} & 1 if compiler conforms to C standard & \texttt{1} \\
\hline
\texttt{\_\_STDC\_VERSION\_\_} & C standard version (as long) & \texttt{201710L} (C17) \\
\hline
\end{tabular}

\subsection{Compiler-Specific Macros}

\textbf{GCC:}
\begin{lstlisting}[language=C]
#ifdef __GNUC__
    #define GCC_VERSION (__GNUC__ * 10000 + __GNUC_MINOR__ * 100 + __GNUC_PATCHLEVEL__)
    #if GCC_VERSION >= 90100
        // GCC 9.1+ specific code
    #endif
#endif
\end{lstlisting}

\textbf{Clang:}
\begin{lstlisting}[language=C]
#ifdef __clang__
    #define CLANG_VERSION (__clang_major__ * 10000 + __clang_minor__ * 100 + __clang_patchlevel__)
#endif
\end{lstlisting}

\textbf{MSVC:}
\begin{lstlisting}[language=C]
#ifdef _MSC_VER
    #if _MSC_VER >= 1920
        // Visual Studio 2019+
    #endif
#endif
\end{lstlisting}

\subsection{Platform Detection Macros}

\begin{lstlisting}[language=C]
// Architecture
#if defined(__x86_64__) || defined(_M_X64)
    #define ARCH_X64
#elif defined(__i386__) || defined(_M_IX86)
    #define ARCH_X86
#elif defined(__arm__) || defined(_M_ARM)
    #define ARCH_ARM
#elif defined(__aarch64__)
    #define ARCH_ARM64
#endif

// Operating System
#if defined(_WIN32) || defined(_WIN64)
    #define OS_WINDOWS
#elif defined(__linux__)
    #define OS_LINUX
#elif defined(__APPLE__)
    #define OS_MACOS
#elif defined(__FreeBSD__)
    #define OS_FREEBSD
#endif
\end{lstlisting}

\subsection{Practical Example: Debug Logging}

\begin{lstlisting}[language=C]
#define LOG_DEBUG(fmt, ...) \
    fprintf(stderr, "[%s:%d] " fmt "\n", __FILE__, __LINE__, ##__VA_ARGS__)

#define LOG_FUNC_ENTRY() \
    fprintf(stderr, "Entering %s() at %s:%d\n", __func__, __FILE__, __LINE__)

void process_data(int value) {
    LOG_FUNC_ENTRY();
    LOG_DEBUG("Processing value: %d", value);
    // ... processing code ...
}

int main(void) {
    process_data(42);
    return 0;
}
\end{lstlisting}

Output:
\begin{verbatim}
Entering process_data() at source.c:10
[source.c:11] Processing value: 42
\end{verbatim}

\section{\#define vs. enum vs. const}

When should you use macros versus other constants?

\subsection{Macros}

\begin{lstlisting}[language=C]
#define MAX_ITEMS 100
#define VERSION "1.0.0"
\end{lstlisting}

\textbf{Pros:}
\begin{itemize}
    \item Works in preprocessor conditionals
    \item No memory allocation (text substitution)
    \item Can be undefined and redefined
\end{itemize}

\textbf{Cons:}
\begin{itemize}
    \item No type checking
    \item Not visible in debugger
    \item Can cause unexpected substitutions
\end{itemize}

\subsection{enum}

\begin{lstlisting}[language=C]
enum {
    MAX_ITEMS = 100,
    BUFFER_SIZE = 1024
};

// Typed enums (C23 or earlier with compiler support)
typedef enum Status {
    STATUS_OK = 0,
    STATUS_ERROR = 1,
    STATUS_PENDING = 2
} Status;
\end{lstlisting}

\textbf{Pros:}
\begin{itemize}
    \item Type-safe (with typed enums)
    \item Visible in debugger
    \item Scoped (with \texttt{enum class} in C++)
\end{itemize}

\textbf{Cons:}
\begin{itemize}
    \item Integer values only
    \item Cannot be used in preprocessor conditionals
\end{itemize}

\subsection{const}

\begin{lstlisting}[language=C]
const int max_items = 100;
const char* const version = "1.0.0";
\end{lstlisting}

\textbf{Pros:}
\begin{itemize}
    \item Type-safe
    \item Scoped properly
    \item Visible in debugger
    \item Can take address
\end{itemize}

\textbf{Cons:}
\begin{itemize}
    \item Cannot be used in preprocessor conditionals
    \item May occupy memory (though compilers often optimize)
    \item Cannot be used for array sizes in C89 (VLA in C99)
\end{itemize}

\subsection{Recommendation}

\begin{lstlisting}[language=C]
// For preprocessor conditionals
#define FEATURE_ENABLED 1
#define PLATFORM_LINUX 1

// For integer constants
enum { MAX_ITEMS = 100, BUFFER_SIZE = 1024 };

// For typed constants
static const char* const APP_NAME = "MyApp";
static const double PI = 3.141592653589793;
\end{lstlisting}

\section{Translation Units}

A \textbf{translation unit} is the result of preprocessing a single source file with all its included headers. This is what the compiler proper processes.

\subsection{Translation Unit Boundaries}

Each \texttt{.c} file is a separate translation unit:

\begin{verbatim}
┌─────────────────┐     ┌─────────────────┐     ┌─────────────────┐
│     main.c      │     │    utils.c      │     │    config.c     │
└────────┬────────┘     └────────┬────────┘     └────────┬────────┘
         │                       │                       │
         ▼                       ▼                       ▼
┌─────────────────┐     ┌─────────────────┐     ┌─────────────────┐
│ Translation     │     │ Translation     │     │ Translation     │
│ Unit 1          │     │ Unit 2          │     │ Unit 3          │
└─────────────────┘     └─────────────────┘     └─────────────────┘
         │                       │                       │
         ▼                       ▼                       ▼
┌─────────────────┐     ┌─────────────────┐     ┌─────────────────┐
│    main.o       │     │    utils.o      │     │    config.o     │
└─────────────────┘     └─────────────────┘     └─────────────────┘
         │                       │                       │
         └───────────────────────┼───────────────────────┘
                                 ▼
                       ┌─────────────────┐
                       │   Executable    │
                       └─────────────────┘
\end{verbatim}

\subsection{Static Variables and Translation Units}

Variables and functions with \texttt{static} linkage are private to their translation unit:

\begin{lstlisting}[language=C]
// utils.c
static int internal_counter = 0;  // Only visible in utils.c

static void internal_helper(void) {  // Only visible in utils.c
    internal_counter++;
}

void public_function(void) {
    internal_helper();
}

// main.c
extern int internal_counter;  // Error! Not accessible
\end{lstlisting}

\subsection{One Definition Rule (ODR)}

Each translation unit can have at most one definition of each external symbol. Violations cause linker errors:

\begin{lstlisting}[language=C]
// header.h
int global_var = 42;  // WRONG! Multiple definitions if included in multiple .c files

// Correct approach:
// header.h
extern int global_var;  // Declaration only

// source.c
int global_var = 42;  // Definition in one place
\end{lstlisting}

\section{Common Preprocessor Idioms}

\subsection{Include Guard with \#pragma once}

\begin{lstlisting}[language=C]
#ifndef MYHEADER_H
#define MYHEADER_H
#pragma once

// Header contents

#endif /* MYHEADER_H */
\end{lstlisting}

This provides both compatibility and optimization.

\subsection{Compile-Time Assertions}

\begin{lstlisting}[language=C]
// C11 _Static_assert
_Static_assert(sizeof(int) == 4, "int must be 4 bytes");
static_assert(sizeof(void*) == 8, "64-bit platform required");

// Pre-C11 hack
#define STATIC_ASSERT(cond, msg) \
    typedef char static_assert_##msg[(cond) ? 1 : -1]

STATIC_ASSERT(sizeof(int) == 4, int_must_be_4_bytes);
\end{lstlisting}

\subsection{X-Macro Pattern}

Generate repetitive code from a data table:

\begin{lstlisting}[language=C]
// Define the list
#define ITEM_LIST \
    X(ITEM_A, 1) \
    X(ITEM_B, 2) \
    X(ITEM_C, 3)

// Generate enum
typedef enum {
    #define X(name, value) name = value,
    ITEM_LIST
    #undef X
} Items;

// Generate string array
const char* item_names[] = {
    #define X(name, value) #name,
    ITEM_LIST
    #undef X
};

// Generate handler functions
#define X(name, value) void handle_##name(void);
ITEM_LIST
#undef X

// Expands to:
// typedef enum { ITEM_A = 1, ITEM_B = 2, ITEM_C = 3 } Items;
// const char* item_names[] = { "ITEM_A", "ITEM_B", "ITEM_C" };
// void handle_ITEM_A(void);
// void handle_ITEM_B(void);
// void handle_ITEM_C(void);
\end{lstlisting}

\subsection{Counting Arguments}

\begin{lstlisting}[language=C]
#define VA_NARGS_IMPL(_1, _2, _3, _4, N, ...) N
#define VA_NARGS(...) VA_NARGS_IMPL(__VA_ARGS__, 4, 3, 2, 1)

VA_NARGS(a)        // 1
VA_NARGS(a, b)     // 2
VA_NARGS(a, b, c)  // 3
\end{lstlisting}

\section{Key Takeaways}

\begin{enumerate}
    \item \textbf{Preprocessing is text-only}: The preprocessor doesn't understand C syntax; it performs pure text substitutions.
    \item \textbf{Parenthesize macro arguments}: Without parentheses, operator precedence can cause unexpected results.
    \item \textbf{Beware of multiple evaluation}: Function-like macros can evaluate arguments multiple times, causing side effects.
    \item \textbf{Use \texttt{\#pragma once} for headers}: It's simpler and less error-prone than traditional include guards.
    \item \textbf{Conditional compilation is powerful}: Use it for platform detection, debug builds, and feature flags.
    \item \textbf{Predefined macros are useful}: \texttt{\_\_FILE\_\_}, \texttt{\_\_LINE\_\_}, and \texttt{\_\_func\_\_} enable powerful debugging macros.
    \item \textbf{Understand translation units}: Each \texttt{.c} file is a separate translation unit, and static symbols are private to their unit.
    \item \textbf{Prefer alternatives when appropriate}: Use \texttt{enum} or \texttt{const} instead of macros for values that don't need preprocessor conditions.
\end{enumerate}

\section{Looking Ahead}

Now that we understand how preprocessing transforms source code into translation units, we're ready to explore what happens next: compilation. In Chapter 3, we'll examine how the compiler parses C code, builds intermediate representations, performs optimizations, and generates assembly code.

\section{Further Reading}

\begin{itemize}
    \item ``The C Programming Language'' (K\&R) - Chapter 4
    \item C17 Standard (ISO/IEC 9899:2018) - Section 6.10 (Preprocessing directives)
    \item GCC Preprocessor Documentation: \url{https://gcc.gnu.org/onlinedocs/cpp/}
    \item ``Modern C'' by Jens Gustedt - Chapter 5 (Preprocessing)
\end{itemize}
